\begingroup
\small
\setlength{\parindent}{0pt}
\setlength{\parskip}{0.55\baselineskip}
\setlength{\abovedisplayskip}{8pt plus 2pt minus 2pt}
\setlength{\belowdisplayskip}{8pt plus 2pt minus 2pt}
\setlength{\jot}{4pt}

\section*{How the Proof Works}

The mechanism driving the proof is visible in one calculation.  Set
\[
V_r:=\partial_t^r u,
\qquad
Q_r:=\partial_t^r p,
\qquad
\Lambda:=(-\Delta)^{1/2}.
\]
Differentiating Navier--Stokes \(r\) times gives
\[
\partial_tV_r
+\sum_{a=0}^{r}\binom ra(V_a\!\cdot\nabla)V_{r-a}
+\nabla Q_r
=\nu\Delta V_r.
\]
At spatial Sobolev height \(s\), define
\[
E_{r,s}(t):=\frac12\|\Lambda^sV_r(t)\|_2^2.
\]
Apply \(\Lambda^s\) to the differentiated equation and pair it with
\(\Lambda^sV_r\).  The time term and the viscous term are
\[
\left\langle\Lambda^s\partial_tV_r,\Lambda^sV_r\right\rangle
=\partial_tE_{r,s},
\qquad
\nu\left\langle\Lambda^s\Delta V_r,\Lambda^sV_r\right\rangle
=-2\nu E_{r,s+1}.
\]
Keep the nonlinear and pressure terms together as
\[
\mathcal P_{r,s}
:=-\operatorname{Re}
\left\langle
\Lambda^sV_r,
\Lambda^s\!\left[
\sum_{a=0}^{r}\binom ra(V_a\!\cdot\nabla)V_{r-a}
+\nabla Q_r
\right]
\right\rangle.
\]
The pairing gives the exact identity
\[
\boxed{
\partial_tE_{r,s}=\mathcal P_{r,s}-2\nu E_{r,s+1},
\qquad
E_{r,s+1}
=\frac{\mathcal P_{r,s}-\partial_tE_{r,s}}{2\nu}.
}
\]
This is the central move.  The time variation of the energy at height \(s\)
contains the energy at height \(s+1\), and viscosity supplies the coefficient
that connects them.  Differentiating again repeats the same move; after \(k\)
steps,
\[
\boxed{
\partial_t^kE_{r,s}
=(-2\nu)^kE_{r,s+k}
+\sum_{j=0}^{k-1}(-2\nu)^j
\partial_t^{\,k-1-j}\mathcal P_{r,s+j}.
}
\]
At the velocity level the same coupling is visible directly:
\[
V_{r+1}
=\nu\Delta V_r
-\sum_{a=0}^{r}\binom ra(V_a\!\cdot\nabla)V_{r-a}
-\nabla Q_r,
\qquad
\nu\Delta:H^{m+1}\longrightarrow H^{m-1}.
\]
The Laplacian is where the temporal and spatial counts meet.  In the velocity
equation, the viscous part of \(V_{r+1}\) contains two spatial derivatives of
\(V_r\).  In the energy identity, integration by parts splits those
derivatives across the two factors and produces \(E_{r,s+1}\).  Repeating the
time differentiation repeats this spatial ascent, while \(\mathcal P_{r,s}\)
carries the complete nonlinear and pressure transfer.

Fix \(N\) and \(M\).  Iterating the displayed recurrence only to those orders
involves finitely many velocity rows, pressure rows, energies, and transfer
terms.  Theorem~\ref{thm:datum-rectangles} carries exactly that finite coupled
system across \((0,T_*)\):
\[
V_n\in L^2(0,T_*;H^{M+1}),
\qquad
V_{n+1}\in L^2(0,T_*;H^{M-1})
\qquad(0\le n\le N).
\]
Finite summation and Proposition~\ref{prop:datum-terminal-realization} give
\[
\mathfrak R^*_{N,M}[u_0,\nu,T_*]
=\sum_{n=0}^{N}
\left(
\|V_n\|_{L^2_tH^{M+1}_x}^2
+\|V_{n+1}\|_{L^2_tH^{M-1}_x}^2
\right)<\infty.
\]
The pressure recurrence supplies the matching \(Q_n\)-coordinates, so this
finite calculation is the rectangle \(\mathcal R_{N,M}[u_0;T_*]\).  Enlarging
\((N,M)\) leaves every shared coordinate unchanged.  The rectangles can thus
be taken together, and for every finite \(r,m\),
\[
\mathcal I[u_0]
\in\varprojlim_{N,M}\mathscr X_{N,M}
\quad\Longrightarrow\quad
u\in H^r(0,T_*;H^m(\mathbb R^3)).
\]

\clearpage

Now name the estimate you want.  A finite-order Sobolev-controlled estimate
specifies finitely many derivatives and target norms.  Write the derivatives as
\[
D_j u:=\partial_t^{a_j}\nabla^{b_j}u,
\qquad 1\le j\le J,
\]
where \(J\), every \(a_j\), and every \(b_j\) are finite.  For a requested norm
\(L_t^{p_j}L_x^{q_j}\), with \(1\le p_j\le\infty\) and
\(2\le q_j\le\infty\), choose finite Sobolev indices \(n_j,m_j\) so that
\[
n_j-a_j>\frac12-\frac1{p_j},
\qquad
m_j-b_j>3\!\left(\frac12-\frac1{q_j}\right).
\]
The time and space Sobolev embeddings then give
\[
H^{n_j}(0,T_*;H^{m_j}(\mathbb R^3))
\hookrightarrow
W^{a_j,p_j}(0,T_*;W^{b_j,q_j}(\mathbb R^3)),
\]
and hence
\[
\|D_j u\|_{L_t^{p_j}L_x^{q_j}}
\le
C_j\|u\|_{H_t^{n_j}H_x^{m_j}}<\infty.
\]
Because the requested list is finite, one can choose
\[
N\ge\max_{1\le j\le J}n_j,
\qquad
M\ge\max_{1\le j\le J}m_j,
\]
and a single rectangle \(\mathcal R_{N,M}[u_0;T_*]\) contains every coordinate
used by the estimate.  Nothing is assembled from unrelated bounds: the
coordinates already agree inside the compatible family.

If the requested expression contains products, choose \(M\) above the
Sobolev algebra threshold and apply the standard product or Moser estimate.
If it contains pressure, use
\[
\begin{aligned}
Q_r
&=R_iR_j\sum_{a=0}^{r}\binom ra V_{a,i}V_{r-a,j},\\
\|Q_r\|_{H^m}
&\lesssim_{r,m}
\sum_{a=0}^{r}\|V_aV_{r-a}\|_{H^m}
\lesssim_{r,m}
\sum_{a=0}^{r}\|V_a\|_{H^m}\|V_{r-a}\|_{H^m},
\qquad m>\frac32.
\end{aligned}
\]
Here the second line uses the Sobolev boundedness of the Riesz transforms.
Spatial Sobolev embedding can then be applied if the requested pressure norm
is \(L^\infty\).  If the estimate asks for an intermediate exponent,
interpolate between two coordinates already present in the rectangle.  If it
asks for time integrability on \((0,T_*)\), use the time Sobolev embedding or
finite-interval H\"older after selecting the required temporal order.

Thus every finite-order Sobolev-controlled request has the same form:
\[
\boxed{
\begin{gathered}
\text{finite derivative and norm requirements}
\longrightarrow
\text{one sufficiently large }(N,M)\\
\longrightarrow
\mathcal R_{N,M}[u_0;T_*]
\longrightarrow
\text{the requested estimate}.
\end{gathered}
}
\]
The requested quantity may be a pointwise derivative bound, a spacetime norm,
a nonlinear product bound, a pressure bound, or a finite collection of them.
The chosen \((N,M)\) and its constant may change with the question.  No
estimate uniform over all derivative orders is being claimed or needed.

For the terminal \(H^m\) trace, select temporal order one and that fixed
spatial height:
\[
u\in H^1(0,T_*;H^m)
\hookrightarrow C([0,T_*];H^m).
\]
Thus \(u(T_*)\in H^m\) for every finite \(m\).  Compatibility identifies
these traces as one endpoint
\[
u_*:=u(T_*)\in\bigcap_{m<\infty}H^m=H^\infty.
\]
Local well-posedness continues the solution from \(u_*\) past \(T_*\), which
contradicts finite maximality and gives \(T_*=\infty\).

The intersection supplies every finite Sobolev window; the estimate being
requested determines which window is read from it.  Its constants may change
with the question, but the solution and the direction of the proof do not.

\endgroup
